% !TeX program = pdflatex
% =====================================================================
%  ISIGK Robotics Club — Management Platform
%  Technical Project Report (Sprint/Chapter based)
%  Compile with: pdflatex rapport_robotic_club.tex (run 2x for TOC)
% =====================================================================
\documentclass[12pt,a4paper]{report}

% ---------------- Packages ----------------
\usepackage[utf8]{inputenc}
\usepackage[T1]{fontenc}
\usepackage[english]{babel}
\usepackage{geometry}
\geometry{margin=2.5cm}
\usepackage{graphicx}
\usepackage{booktabs}
\usepackage{longtable}
\usepackage{array}
\usepackage{multirow}
\usepackage{xcolor}
\usepackage{enumitem}
\usepackage{fancyhdr}
\usepackage{titlesec}
\usepackage{caption}
\usepackage{tikz}
\usetikzlibrary{shapes,arrows,positioning}
\usepackage[hidelinks]{hyperref}

\graphicspath{{./figures/}}

% ---------------- Custom commands ----------------
\newcommand{\todo}[1]{\textcolor{red}{\textbf{[TODO: #1]}}}
\newcommand{\vprio}{\textcolor{red}{\textbf{MVP}}}
\newcommand{\pprio}{\textcolor{blue}{\textbf{PROD}}}
\newcommand{\testprio}{\textcolor{orange}{\textbf{TEST}}}

% ---------------- Headers / footers ----------------
\pagestyle{fancy}
\fancyhf{}
\fancyhead[L]{\small ISIGK Robotics Club — Management Platform}
\fancyhead[R]{\small \thepage}
\renewcommand{\headrulewidth}{0.4pt}

\hypersetup{
    pdftitle={ISIGK Robotics Club Management Platform},
    pdfauthor={Project Team},
    colorlinks=true,
    linkcolor=black!70!blue,
    urlcolor=black!60!blue
}

\begin{document}

% =====================================================================
%  TITLE PAGE
% =====================================================================
\begin{titlepage}
    \centering
    \vspace*{1cm}

    {\Large \textbf{Institut Supérieur d'Informatique et de Gestion (ISIGK)}\par}
    \vspace{0.5cm}
    {\large Department of Computer Science\par}

    \vspace{2cm}
    {\Huge \textbf{Robotics Club\\[0.2cm] Management Platform}\par}
    \vspace{1cm}
    {\Large Technical Project Report\par}

    \vspace{2cm}
    {\large
    \begin{tabular}{rl}
        Project Supervisors : & Pr. [Name] \& [Name] \\
        Club Leader         : & [Name] \\
        Co-Leader           : & [Name] \\
        Development Team    & [Names] \\
        Academic Year       : & 2026 -- 2027 \\
    \end{tabular}\par}

    \vfill
    {\large \today\par}
\end{titlepage}

% =====================================================================
%  TABLE OF CONTENTS
% =====================================================================
\tableofcontents
\listoftables
\listoffigures

\clearpage

% =====================================================================
%  CHAPTER 1 — GENERAL INTRODUCTION
% =====================================================================
\chapter{General Introduction}

\section{Project Overview}

The Robotics Club of ISIGK is a student association whose mission is to promote robotics,
electronics and innovation among students through workshops, competitions and collaborative
projects. The club is governed by a well-defined hierarchical organization: one professor
acting as an external supervisor, one leader and one co-leader directing the strategy, and
five managers each in charge of a specific department (Human Resources, Events, Media,
Finance and Materials). Each department contains members, and the club also welcomes guest
members who are in an observation or trial period.

This project aims to design and develop a \textbf{web-based management platform} that
digitizes the club's daily operations. The platform provides each actor with a dedicated
interface adapted to his role and responsibilities: member registry management, event
organization, internal communication through chat, budget and inventory tracking, media
publication, and real-time notifications.

The deliverable of this project follows an \textbf{agile, sprint-based approach}: each
sprint is documented as a chapter of this report, allowing continuous tracking of the
progress from the initial backlog to the final production-ready product.

\section{Problem Statement}

Despite its dynamic organization, the Robotics Club currently suffers from a fragmented
and informal management style:

\begin{itemize}[leftmargin=*]
    \item \textbf{Scattered communication:} exchanges take place in multiple WhatsApp and
          Telegram groups, mixing administrative, technical and social content. Important
          information gets lost and guests or new members are often excluded.
    \item \textbf{No structured information system:} member lists, attendance and
          department data are stored in personal files (Excel, Word) with no single source
          of truth and no access control.
    \item \textbf{Manual event organization:} invitations, RSVP tracking and post-event
          follow-up are done manually, leading to duplicated effort and low participation
          visibility.
    \item \textbf{Untracked resources:} materials, electronic components and tools are
          borrowed without traceability, causing losses and stock discrepancies.
    \item \textbf{Opaque finances:} income, expenses and per-department budgets are not
          visible to the club's supervision, limiting accountability.
    \item \textbf{No remote collaboration:} there is no integrated solution for video
          meetings tied to the club's departments and events.
\end{itemize}

Existing general-purpose tools (WhatsApp, Google Drive, Discord, Notion) each cover only a
fragment of these needs and do not enforce the club's specific role hierarchy. There is
therefore a real need for a \textbf{custom platform aligned with the club's organization}.

\section{Project Objectives}

The main objective is to build a single platform that makes it easy for every user to
manage his role within the club. The specific objectives are:

\begin{enumerate}[leftmargin=*]
    \item \textbf{O1 -- Role-based access:} implement a strict permission model
          (Professor, Leader, Co-leader, Manager, Member, Guest) so that each user sees and
          does exactly what his role allows.
    \item \textbf{O2 -- Department management:} provide the five managers with the tools
          to run their departments (HR, Events, Media, Finance, Materials).
    \item \textbf{O3 -- Unified communication:} integrate internal chat (direct messages
          and channels) with file and image sharing.
    \item \textbf{O4 -- Collaboration:} enable video meetings, tested first with an
          open-source solution before the final choice.
    \item \textbf{O5 -- Real-time awareness:} deliver in-app notifications for messages,
          events, tasks and alerts.
    \item \textbf{O6 -- Guest integration:} offer a smooth pipeline from public visitor
          to guest, then to confirmed club member.
    \item \textbf{O7 -- Traceability:} keep history of activities, finances and material
          loans for accountability towards the supervisor.
    \item \textbf{O8 -- Progressive delivery:} ship a Minimum Viable Product (MVP) first,
          then iterate towards a production-grade version (PROD).
\end{enumerate}

\clearpage

% =====================================================================
%  CHAPTER 2 — PROJECT CONTEXT
% =====================================================================
\chapter{Project Context}

\section{Project Scope and Constraints}

\subsection{Scope}

The platform covers the following functional perimeter:

\begin{center}
\begin{tabular}{>{\bfseries}p{3.5cm} p{9.5cm}}
\toprule
\textbf{In scope} & \textbf{Out of scope (v1)} \\
\midrule
Authentication & Payment gateway integration \\
Role management (6 roles) & Native mobile application \\
Departments \& members registry & Email/Telegram notifications (PROD) \\
Guest application pipeline & Advanced analytics / BI \\
Events \& RSVP & Self-hosted video infrastructure \\
Tasks assignment & Full-text search engine \\
Chat (DM, channels, attachments) & Multi-language UI (PROD) \\
In-app notifications & \\
Finance \& materials tracking & \\
Media gallery \& announcements & \\
Public landing page & \\
\bottomrule
\end{tabular}
\end{center}

\subsection{Constraints}

\begin{itemize}[leftmargin=*]
    \item \textbf{Human:} a small student development team, working alongside academic
          obligations; sprints must be realistic (1--2 weeks).
    \item \textbf{Technical:} video conferencing must not be self-hosted in v1 (bandwidth
          and maintenance costs); an external open-source service (Jitsi) is evaluated
          first.
    \item \textbf{Financial:} the budget is limited to free tiers and open-source
          tooling.
    \item \textbf{Organizational:} the platform must respect the club hierarchy; the
          professor has read-only access to everything, including chats.
    \item \textbf{Legal:} personal data of members must be handled responsibly
          (password hashing, restricted access).
\end{itemize}

\subsection{Assumptions}

\begin{itemize}[leftmargin=*]
    \item The club structure (1 professor, 1 leader, 1 co-leader, 5 managers, members,
          guests) is stable for the duration of the project.
    \item The user base remains in the order of tens of users, not thousands.
    \item Users access the platform mainly from laptops and smartphones via a web
          browser.
\end{itemize}

\section{Chapter Conclusion}

The project is bounded by tight human and financial constraints but targets a clearly
identified pain: the absence of a unified, role-aware information system for the club.
By restricting the scope to a web application with proven open-source technologies, and
by deferring heavy features (self-hosted video, email alerts) to the production phase,
the project remains achievable by a student team while delivering real value in the MVP.

\clearpage

% =====================================================================
%  CHAPTER 3 — STATE OF THE ART
% =====================================================================
\chapter{State of the Art}

\section{Existing solutions}

Several categories of existing solutions were analyzed:

\subsection{General messaging tools}
\textbf{WhatsApp}, \textbf{Telegram} and \textbf{Discord}~\cite{discord} dominate student
association communication. They provide groups, media sharing and voice channels, but
offer no role hierarchy matching a club organization, no member registry, and no
traceability of administrative data. Discord comes closest with text/voice channels and
role tags, yet it remains a communication tool and not a management system.

\subsection{Collaboration and productivity suites}
\textbf{Google Workspace} (Drive, Sheets, Forms) and \textbf{Notion} allow clubs to store
documents, build simple databases and forms~\cite{notion}. They are flexible but require
manual configuration, lack a coherent permission model, and provide no integrated chat,
notifications or event workflow.

\subsection{Project and task management}
\textbf{Trello}, \textbf{Jira} and \textbf{Asana}~\cite{trello} offer boards, tasks and
deadlines. They fit development teams better than club organizations, and again provide
no finance, inventory or membership features.

\subsection{Club / association management software}
Products such as \textbf{WildApricot}, \textbf{MembershipWorks} or management modules of
\textbf{Moodle}~\cite{moodle} address membership, dues and events. They are designed for
Western associations, are often paid per member, and do not natively support department
hierarchies, chat or inventory of electronic components.

\subsection{Communication infrastructure technologies}
For real-time features, the state of the art converges on \textbf{WebSocket}
technologies, commonly implemented with \textbf{Socket.io} on Node.js
applications~\cite{socketio}. For video meetings, the dominant open-source option is
\textbf{Jitsi Meet}~\cite{jitsi}, which can be embedded in a web page through an iframe
without operating any media server, at zero cost.

\section{Synthesis and Positioning of the project}

\begin{center}
\small
\begin{tabular}{l c c c c c c}
\toprule
\textbf{Solution} & \textbf{Roles} & \textbf{Chat} & \textbf{Events} &
\textbf{Finance} & \textbf{Stock} & \textbf{Free} \\
\midrule
WhatsApp / Telegram & -- & \checkmark & partial & -- & -- & \checkmark \\
Discord             & partial & \checkmark & -- & -- & -- & \checkmark \\
Google Drive/Sheets & -- & -- & partial & partial & partial & \checkmark \\
Notion              & partial & -- & partial & partial & partial & partial \\
Trello / Jira       & partial & -- & -- & -- & -- & partial \\
WildApricot / Moodle& \checkmark & -- & \checkmark & \checkmark & -- & -- \\
\textbf{Our platform} & \checkmark & \checkmark & \checkmark & \checkmark &
\checkmark & \checkmark \\
\bottomrule
\end{tabular}
\end{center}

No single existing solution combines: (i) the exact club hierarchy with a strict
read-only supervisor, (ii) department-scoped management modules, (iii) chat and
notifications, and (iv) zero operating cost. The project therefore positions itself as a
\textbf{lightweight, open-source, role-aware club information system}, built with a
student-friendly JavaScript stack and leveraging external free services (Jitsi) for heavy
features.

\section{Wrap-up}

The state of the art confirms the feasibility of every targeted feature: real-time chat
with Socket.io, embedded video with Jitsi, and role-based web applications with
component-based front-end frameworks. The differentiation of this project lies in its
\textbf{integration} and in its \textbf{permission model}, not in technological
novelty. The next chapter formalizes these needs into requirements, a product backlog and
a sprint plan.

\clearpage

% =====================================================================
%  CHAPTER 4 — REQUIREMENTS & PLANNING
% =====================================================================
\chapter{Requirements \& Planning}

\section{Methodology}

The project follows \textbf{Scrum}~\cite{scrumguide} with short sprints:

\begin{itemize}[leftmargin=*]
    \item \textbf{Product Owner:} the club Leader (validates the backlog and priorities).
    \item \textbf{Scrum Master:} the development team lead (facilitates the process).
    \item \textbf{Development Team:} the student developers.
    \item \textbf{Sprint duration:} 1 to 2 weeks; each sprint is reported as one
          ``Sprint/Chapter'' of this document.
    \item \textbf{Artifacts:} product backlog (Section~\ref{sec:backlog}), sprint backlog,
          increment (working software).
    \item \textbf{Ceremonies:} sprint planning, daily sync, sprint review with the club
          leaders, sprint retrospective (Section 5.5 of each chapter).
    \item \textbf{Definition of Done:} feature coded, tested manually, documented and
          demoed to the Product Owner.
\end{itemize}

Priorities use the MoSCoW-like notation: \vprio{} (must have in the MVP),
\testprio{} (evaluated experimentally), \pprio{} (deferred to the production version).

\section{Stakeholder Identification}

\begin{center}
\begin{tabular}{>{\bfseries}p{2.6cm} p{3.2cm} p{7.4cm}}
\toprule
\textbf{Stakeholder} & \textbf{Role w.r.t. project} & \textbf{Involvement} \\
\midrule
Professor        & Supervisor / end user (read-only) & Validates that the club
    activities are traceable and visible \\
Club Leader      & Product Owner & Prioritizes the backlog, validates increments \\
Co-Leader        & Deputy Product Owner & Covers the Leader's validations \\
HR Manager       & End user & Manages members, attendance, applications \\
Events Manager   & End user & Manages events, RSVPs, reminders \\
Media Manager    & End user & Manages the gallery and publications \\
Finance Manager  & End user & Manages income, expenses, budgets \\
Materials Manager& End user & Manages inventory and loans \\
Members          & End users & Use tasks, chats, events \\
Guests           & External users & Public site, application form \\
Development Team & Builders & Designs, develops, tests and deploys the platform \\
ISIGK            & Host institution & Provides identity and context \\
\bottomrule
\end{tabular}
\end{center}

\section{Functional Requirements}

\begin{longtable}{>{\bfseries}p{1.4cm} p{9.2cm} p{2.4cm} p{1.6cm}}
\toprule
\textbf{ID} & \textbf{Requirement} & \textbf{Priority} & \textbf{Sprint} \\
\midrule
\endfirsthead
\toprule
\textbf{ID} & \textbf{Requirement} & \textbf{Priority} & \textbf{Sprint} \\
\midrule
\endhead
FR-01 & Register / login with email and password & \vprio & S1 \\
FR-02 & Assign, promote and demote roles (guest $\rightarrow$ member $\rightarrow$ manager) & \vprio & S1 \\
FR-03 & Professor account with global \emph{read-only} access (data + chats) & \vprio & S1 \\
FR-04 & Manage the 5 departments and the Board team & \vprio & S1 \\
FR-05 & Members directory with search and department filter & \vprio & S2 \\
FR-06 & User profile (name, role, department, photo, bio) & \vprio & S2 \\
FR-07 & Guest application form and approval pipeline & \vprio & S2 \\
FR-08 & Create / edit / delete events & \vprio & S3 \\
FR-09 & RSVP and attendee list per event & \vprio & S3 \\
FR-10 & Assign tasks with deadline and status & \vprio & S3 \\
FR-11 & Direct messages between users & \vprio & S4 \\
FR-12 & Department channels (auto-joined) and club-wide channel & \vprio & S4 \\
FR-13 & Private Board channel (Leader, Co-leader, Professor) & \vprio & S4 \\
FR-14 & File and image attachments inside chat & \vprio & S4 \\
FR-15 & Real-time in-app notification center & \vprio & S3--S4 \\
FR-16 & Record income and expenses with proof & \vprio & S5 \\
FR-17 & Per-department budgets and monthly report & \vprio & S5 \\
FR-18 & Materials inventory with quantities and states & \vprio & S5 \\
FR-19 & Loan / reservation of materials with return date & \vprio & S5 \\
FR-20 & Low-stock automatic alert & \vprio & S5 \\
FR-21 & Media gallery (photos, videos) & \vprio & S6 \\
FR-22 & Public landing page and join form & \vprio & S6 \\
FR-23 & Announcements from Leader broadcast to all & \vprio & S6 \\
FR-24 & Role-adapted dashboards (KPIs) & \vprio & S2--S6 \\
FR-25 & Video meeting button per channel / event & \testprio & S4 (test) \\
FR-26 & Email notifications & \pprio & PROD \\
FR-27 & Mentions and read receipts & \pprio & PROD \\
FR-28 & Kanban task board & \pprio & PROD \\
FR-29 & Member dues tracking & \pprio & PROD \\
FR-30 & Purchase requests (Materials $\rightarrow$ Finance) & \pprio & PROD \\
FR-31 & Extra sub-channels created by managers & \pprio & PROD \\
FR-32 & Post-event reports with photos & \pprio & PROD \\
FR-33 & Social media posts planner & \pprio & PROD \\
FR-34 & Multi-language interface (FR / EN) & \pprio & PROD \\
FR-35 & Notification preferences per user & \pprio & PROD \\
\bottomrule
\end{longtable}

\section{Non Functional Requirements}

\begin{center}
\begin{tabular}{>{\bfseries}p{2.8cm} p{11cm}}
\toprule
\textbf{Category} & \textbf{Requirement} \\
\midrule
Security & Passwords hashed (bcrypt), JWT sessions, strict server-side role checks,
           protection against SQL injection and XSS \\
Confidentiality & A user only accesses his department's data; the Leader, Co-leader and
           Professor access everything; the Professor cannot modify anything \\
Performance & Pages load in $<$ 2s on a standard campus connection; chat messages
           delivered in $<$ 500ms \\
Availability & Hosted on a free-tier service with $\geq$ 95\% availability; graceful
           degradation if the video service is down \\
Usability & Responsive interface (desktop + mobile), French/English friendly,
           discoverable navigation per role \\
Reliability & Daily database backups; no data loss on server restart \\
Scalability & Architecture able to serve 100+ users without redesign \\
Maintainability & Modular code (one module per department), documented API, version
           control with Git \\
Portability & Runs on any system with a modern browser; deployable on Linux hosting \\
\bottomrule
\end{tabular}
\end{center}

\section{Product Backlog}
\label{sec:backlog}

\subsection{Epic summary}

\begin{center}
\begin{tabular}{>{\bfseries}p{1.2cm} p{4.6cm} p{6.4cm}}
\toprule
\textbf{Epic} & \textbf{Title} & \textbf{Contains} \\
\midrule
E1  & Authentication \& Roles     & Login, roles, professor read-only \\
E2  & Departments \& Members      & Registry, directory, profiles, guest pipeline \\
E3  & Events                      & Events, RSVP, reminders \\
E4  & Tasks                       & Assignment, deadlines, statuses \\
E5  & Chat                        & DM, channels, attachments \\
E6  & Video Meetings              & Jitsi experiment, meeting buttons \\
E7  & Notifications               & In-app center, real-time \\
E8  & Finance                     & Income, expenses, budgets, reports \\
E9  & Materials                   & Inventory, loans, stock alerts \\
E10 & Media                       & Gallery, public preview \\
E11 & Announcements               & Leader broadcasts \\
E12 & Public Site                 & Landing page, join form \\
E13 & Settings                    & Profile, password, preferences \\
E14 & Dashboards                  & Role-adapted KPIs \\
\bottomrule
\end{tabular}
\end{center}

\subsection{User stories — MVP}

\begin{longtable}{>{\bfseries}p{1.5cm} p{9.6cm} p{2.6cm}}
\toprule
\textbf{ID} & \textbf{As a \dots I want to \dots so that \dots} & \textbf{Epic} \\
\midrule
\endfirsthead
\toprule
\textbf{ID} & \textbf{Story} & \textbf{Epic} \\
\midrule
\endhead
US-1.1 & As a \textbf{visitor} I want to register and log in so that I access the platform & E1 \\
US-1.2 & As a \textbf{leader} I want to assign roles so that each user gets the right permissions & E1 \\
US-1.3 & As a \textbf{professor} I want to view all data and chats read-only so that I supervise the club & E1 \\
US-2.1 & As a \textbf{leader} I want the 5 departments created so that the club is structured & E2 \\
US-2.2 & As a \textbf{manager} I want to see the members of my department so that I coordinate them & E2 \\
US-2.3 & As a \textbf{member} I want a profile so that others identify me & E2 \\
US-2.4 & As a \textbf{guest} I want to apply to join a department so that I become a member & E2 \\
US-2.5 & As a \textbf{manager} I want to accept or refuse applications so that I control my team & E2 \\
US-3.1 & As an \textbf{events manager} I want to create events so that the club is informed & E3 \\
US-3.2 & As a \textbf{member} I want to RSVP so that organizers plan properly & E3 \\
US-3.3 & As a \textbf{member} I want to receive event reminders so that I don't miss them & E3 \\
US-4.1 & As a \textbf{manager} I want to assign tasks to my members so that work progresses & E4 \\
US-4.2 & As a \textbf{leader} I want to assign tasks to managers so that objectives are tracked & E4 \\
US-4.3 & As a \textbf{member} I want to update my task status so that everyone sees progress & E4 \\
US-5.1 & As a \textbf{member} I want to send direct messages so that I communicate privately & E5 \\
US-5.2 & As a \textbf{member} I want a department channel so that my team discusses in one place & E5 \\
US-5.3 & As a \textbf{member} I want a club-wide channel so that news reaches everyone & E5 \\
US-5.4 & As a \textbf{leader} I want a private Board channel so that strategy stays confidential & E5 \\
US-5.5 & As a \textbf{member} I want to share files and images in chat so that we exchange documents & E5 \\
US-7.1 & As a \textbf{user} I want a notification center so that I never miss anything & E7 \\
US-7.2 & As a \textbf{user} I want real-time notifications so that I don't refresh the page & E7 \\
US-8.1 & As a \textbf{finance manager} I want to record income and expenses so that finances are traceable & E8 \\
US-8.2 & As a \textbf{leader} I want per-department budgets so that spending is controlled & E8 \\
US-8.3 & As a \textbf{professor} I want a monthly financial report so that I supervise the club budget & E8 \\
US-9.1 & As a \textbf{materials manager} I want an inventory so that stock is known & E9 \\
US-9.2 & As a \textbf{member} I want to borrow material with a return date so that items come back & E9 \\
US-9.3 & As a \textbf{materials manager} I want low-stock alerts so that I reorder in time & E9 \\
US-10.1 & As a \textbf{media manager} I want to upload photos and videos so that activities are showcased & E10 \\
US-10.2 & As a \textbf{visitor} I want to see a public gallery preview so that I discover the club & E10 \\
US-11.1 & As a \textbf{leader} I want to post announcements so that the whole club is informed & E11 \\
US-12.1 & As a \textbf{visitor} I want a landing page so that I understand the club & E12 \\
US-12.2 & As a \textbf{visitor} I want a join form so that I become a guest & E12 \\
US-13.1 & As a \textbf{user} I want to edit my profile and password so that my data stays correct & E13 \\
US-14.1 & As a \textbf{leader} I want club KPIs so that I pilot the club & E14 \\
US-14.2 & As a \textbf{manager} I want department KPIs so that I pilot my team & E14 \\
US-14.3 & As a \textbf{member} I want my tasks and my department news in one place & E14 \\
\bottomrule
\end{longtable}

\subsection{User stories — Video experiment (TEST)}

\begin{center}
\begin{tabular}{>{\bfseries}p{1.5cm} p{10cm} p{2.2cm}}
\toprule
\textbf{ID} & \textbf{Story} & \textbf{Epic} \\
\midrule
US-6.1 & As a \textbf{user} I want a ``Start meeting'' button in a channel so that we talk face to face & E6 \\
US-6.2 & As a \textbf{manager} I want a permanent room for my department so that meetings are easy & E6 \\
US-6.3 & As a \textbf{member} I want a ``Join meeting'' button on events so that I attend remotely & E6 \\
\bottomrule
\end{tabular}
\end{center}

\textbf{Decision gate:} Jitsi Meet is prototyped in Sprint~4; the team accepts or rejects
it in the Sprint~4 review before integrating it into the MVP.

\subsection{User stories — Production version (PROD)}

\begin{center}
\begin{tabular}{>{\bfseries}p{1.5cm} p{10cm} p{2.2cm}}
\toprule
\textbf{ID} & \textbf{Story} & \textbf{Epic} \\
\midrule
US-6.4 & As a \textbf{user} I want meetings integrated natively & E6 \\
US-7.3 & As a \textbf{user} I want email notifications for important events & E7 \\
US-7.4 & As a \textbf{user} I want Telegram alerts so that I stay informed on my phone & E7 \\
US-5.6 & As a \textbf{user} I want @mentions and read receipts & E5 \\
US-5.7 & As a \textbf{manager} I want to create sub-channels in my department & E5 \\
US-4.4 & As a \textbf{manager} I want a kanban board so that I visualize tasks & E4 \\
US-8.4 & As a \textbf{finance manager} I want to track member dues & E8 \\
US-9.4 & As a \textbf{materials manager} I want purchase requests sent to Finance & E9 \\
US-3.4 & As an \textbf{organizer} I want a post-event report with photos & E3 \\
US-10.3 & As a \textbf{media manager} I want a social posts planner & E10 \\
US-12.3 & As a \textbf{visitor} I want the site in French and English & E12 \\
US-13.2 & As a \textbf{user} I want notification preferences & E13 \\
\bottomrule
\end{tabular}
\end{center}

\section{Sprint / Chapter Planning}
\label{sec:sprintplan}

Each sprint is delivered and documented as one chapter of this report.

\begin{longtable}{>{\bfseries}p{1.3cm} p{5.2cm} p{5.6cm} p{1.6cm}}
\toprule
\textbf{Sprint} & \textbf{Goal} & \textbf{Main deliverables} & \textbf{Chapter} \\
\midrule
\endfirsthead
\toprule
\textbf{Sprint} & \textbf{Goal} & \textbf{Main deliverables} & \textbf{Chapter} \\
\midrule
\endhead
S0 & Project setup & Repo, environment, skeleton app, this backlog & Ch. 5 \\
S1 & Access control & Auth, 6 roles, RBAC, profiles & Ch. 6 \\
S2 & Members & Departments, directory, guest pipeline, dashboard base & Ch. 7 \\
S3 & Organization & Events, RSVP, tasks, notifications core & Ch. 8 \\
S4 & Communication & Chat (DM, channels, attachments) + video test (Jitsi) & Ch. 9 \\
S5 & Resources & Finance, budgets, inventory, loans, stock alerts & Ch. 10 \\
S6 & Visibility & Media gallery, announcements, public landing page & Ch. 11 \\
S7 & Hardening & Testing, bug fixing, deployment, demo data, handover & Ch. 12 \\
--- & \textbf{PROD phase} & Email, Telegram, kanban, dues, sub-channels\dots & Ch. 13+ \\
\bottomrule
\end{longtable}

\section{Conclusion}

This chapter transformed the club's needs into a formal requirement set: 35 functional
requirements, 9 categories of non-functional requirements, a backlog of 14 epics and more
than 50 user stories split between MVP, experimental and production scopes, and a
sprint plan aligned with the report structure. The project is now ready to enter Sprint~0,
which initializes the technical environment and the application skeleton.

\clearpage

% =====================================================================
%  CHAPTER 5 — SPRINT / CHAPTER 0
% =====================================================================
\chapter{Sprint / Chapter 0}

\section{Introduction}

Sprint~0 is an initialization sprint: it does not deliver business features but sets up
everything the following sprints depend on -- repository, environment, application
skeleton, common conventions and the product backlog validated with the Product Owner.

\textbf{Sprint goal:} ``A running skeleton application with authentication stubs and a
validated backlog.''

\section{Actors}

The actors of the system are formalized as follows (they are refined in every sprint):

\begin{center}
\begin{tabular}{>{\bfseries}p{2.4cm} p{5.2cm} p{5.8cm}}
\toprule
\textbf{Actor} & \textbf{Description} & \textbf{Main interactions} \\
\midrule
Professor     & Club supervisor, read-only on everything & View dashboards, data, chats \\
Leader        & Club head, full control & Manage all, assign roles, validate \\
Co-Leader     & Deputy head & Same powers as the Leader \\
HR Manager    & Head of Human Resources & Members, attendance, applications \\
Events Manager& Head of Events & Events, RSVP, reminders \\
Media Manager & Head of Media & Gallery, publications \\
Finance Manager & Head of Finance & Income, expenses, budgets \\
Materials Manager & Head of Materials & Inventory, loans, alerts \\
Member        & Department member & Tasks, chats, events, loans \\
Guest         & Visitor / trial user & Public site, application, limited chat \\
\bottomrule
\end{tabular}
\end{center}

\section{Diagrams}

\begin{figure}[htbp]
    \centering
    \fbox{\parbox[c][6cm][c]{0.92\linewidth}{\centering
        \textbf{Use Case Diagram (UML)}\\[0.4cm]
        \todo{Insert \texttt{usecase\_global.png} after Sprint 0 review}\\
        Actors: Professor, Leader, Co-leader, 5 Managers, Member, Guest
    }}
    \caption{Global use case diagram — to be delivered in Sprint 0}
    \label{fig:usecase}
\end{figure}

\begin{figure}[htbp]
    \centering
    \fbox{\parbox[c][6cm][c]{0.92\linewidth}{\centering
        \textbf{Class / Data Model Diagram}\\[0.4cm]
        \todo{Insert entity-relationship diagram (users, departments, events,
        messages, notifications, expenses, materials, \dots)}
    }}
    \caption{Data model — to be delivered in Sprint 0}
    \label{fig:datamodel}
\end{figure}

\begin{figure}[htbp]
    \centering
    \fbox{\parbox[c][5cm][c]{0.92\linewidth}{\centering
        \textbf{Architecture Diagram (3 tiers)}\\[0.4cm]
        \todo{Browser (React) $\rightarrow$ API (Node/Express) $\rightarrow$ Database;
        Socket.io channel; Jitsi iframe}
    }}
    \caption{Technical architecture — to be delivered in Sprint 0}
    \label{fig:arch}
\end{figure}

\section{Technology Stack}

\textbf{To be defined and justified in the Sprint~0 review with the Product Owner.}
The shortlist discussed so far:

\begin{center}
\begin{tabular}{p{3.4cm} p{10cm}}
\toprule
\textbf{Layer} & \textbf{Candidates} \\
\midrule
Front-end   & React + Vite, or Next.js \\
Back-end    & Node.js + Express, or Next.js API routes \\
Database    & SQLite (dev) $\rightarrow$ PostgreSQL (hosted) \\
Real-time   & Socket.io \\
Video       & Jitsi Meet (embedded) — \testprio decision gate \\
Auth        & JWT + bcrypt \\
Deployment  & Free tier (Vercel / Railway / local server) \\
\bottomrule
\end{tabular}
\end{center}

\todo{Fill this table with the final choices and add a justification paragraph after
the Sprint 0 review.}

\section{Review \& Retrospective}

\textbf{Sprint review:} \todo{to be filled after the review meeting with the Product
Owner (date, demo feedback, backlog adjustments).}

\begin{center}
\begin{tabular}{p{4.6cm} p{4.6cm} p{4.6cm}}
\toprule
\textbf{What went well} & \textbf{What went wrong} & \textbf{Actions to improve} \\
\midrule
\todo{} & \todo{} & \todo{} \\
\bottomrule
\end{tabular}
\end{center}

\section{Conclusion}

Sprint~0 \todo{conclusion to be written after the sprint: environment status, backlog
validation, decisions taken (stack, video solution), and outlook for Sprint~1.}

\clearpage

% =====================================================================
%  CHAPTER 6 — SPRINT / CHAPTER 1
% =====================================================================
\chapter{Sprint / Chapter 1}

\section{Introduction}
\todo{Sprint 1 goal: authentication, the 6 roles and role-based access control.}

\section{Actors}
\todo{Actors concerned by this sprint and their permissions matrix.}

\section{Diagrams}
\todo{Use case (auth), sequence diagram (login), class diagram (User, Role).}

\section{Technology Stack}
\todo{Technologies used and difficulties encountered for this sprint.}

\section{Review \& Retrospective}
\todo{Review feedback and retrospective table.}

\section{Conclusion}
\todo{Sprint 1 conclusion and outlook for Sprint 2.}

% =====================================================================
%  (Further sprint chapters are duplicated with the same 6 sections)
% =====================================================================
%  CHAPTER 7  — Sprint 2 : Members, departments, guest pipeline
%  CHAPTER 8  — Sprint 3 : Events, tasks, notifications
%  CHAPTER 9  — Sprint 4 : Chat + video test (Jitsi)
%  CHAPTER 10 — Sprint 5 : Finance + Materials
%  CHAPTER 11 — Sprint 6 : Media, announcements, public site
%  CHAPTER 12 — Sprint 7 : Testing, deployment, handover
%  CHAPTER 13 — Production phase
% =====================================================================

\clearpage

% =====================================================================
%  GENERAL CONCLUSION
% =====================================================================
\chapter{General Conclusion}

\todo{To be written at the end of the project: recap of the initial problem, the
delivered MVP, the PROD outlook, measured results (users, feedback from the club,
metrics), individual learnings and perspectives (mobile app, native video, inter-club
federation).}

\clearpage

% =====================================================================
%  BIBLIOGRAPHY
% =====================================================================
\chapter*{Bibliography}
\addcontentsline{toc}{chapter}{Bibliography}

\begin{enumerate}[label={[\arabic*]}]
    \item K. Schwaber and J. Sutherland, \emph{The Scrum Guide: The Definitive Guide to
          Scrum: The Rules of the Game}, 2020. [Online]. Available:
          \url{https://www.scrumguides.org/scrum-guide.html}
    \item R. S. Pressman and B. R. Maxim, \emph{Software Engineering: A Practitioner's
          Approach}, 9th ed. McGraw-Hill, 2019.
    \item I. Sommerville, \emph{Software Engineering}, 10th ed. Pearson, 2015.
    \item C. Larman, \emph{Applying UML and Patterns: An Introduction to Object-Oriented
          Analysis and Design and Iterative Development}, 3rd ed. Prentice Hall, 2004.
    \item Mozilla Web Docs, ``The WebSocket API.'' [Online]. Available:
          \url{https://developer.mozilla.org/en-US/docs/Web/API/WebSockets_API}
\end{enumerate}

% =====================================================================
%  APPENDICES
% =====================================================================
\clearpage
\appendix
\chapter{Appendix A — Database draft}
\todo{Physical data model, table descriptions and sample SQL schema.}

\chapter{Appendix B — Screenshots}
\todo{Screenshots of each sprint increment.}

\chapter{Appendix C — Meeting minutes}
\todo{Sprint reviews and Product Owner validation minutes.}

\chapter{Appendix D — Backlog export}
\todo{CSV/board export of the product backlog.}

% =====================================================================
%  WEBOGRAPHY
% =====================================================================
\chapter*{Webography}
\addcontentsline{toc}{chapter}{Webography}

\begin{itemize}[leftmargin=*]
    \item Scrum Guide. \url{https://www.scrumguides.org/}
    \item React — official documentation. \url{https://react.dev/}
    \item Vite. \url{https://vitejs.dev/}
    \item Node.js. \url{https://nodejs.org/}
    \item Express.js. \url{https://expressjs.com/}
    \item Socket.io. \url{https://socket.io/}
    \item Jitsi Meet handbook. \url{https://jitsi.github.io/handbook/}
    \item SQLite. \url{https://www.sqlite.org/}
    \item PostgreSQL. \url{https://www.postgresql.org/}
    \item JSON Web Tokens. \url{https://jwt.io/}
    \item bcrypt. \url{https://www.npmjs.com/package/bcrypt}
\end{itemize}

\end{document}
